\section{本章小结}
\label{sec:sgd-summary}

本章回答了一个问题：不读全部目标项时，怎样构造更新并判断它的结果。抽样决定读取哪些项，读数的条件无偏是一切分析的起点，但单步读数可能指向上升方向，要靠期望递推和步长控制来管理噪声。

方差控制决定怎样利用额外读取和记忆。小批量、非均匀抽样和控制变量都能减小读数的方差，代价是快照、梯度表或抽样分布的维护。方差减小并不等同于总计算减少，要按相同的目标和精度比较。

动量、自适应尺度和方向处理分别改变实际更新的不同环节。动量把历史读数合成方向，平滑噪声的同时带来滞后，它对当前梯度不是无偏的，常步长下的噪声残余由有效步长决定。Nesterov方法的读数对应前瞻点，确定性的加速不能直接继承，噪声在势函数中留下不随步数消失的项。Lion与Muon对动量做非线性变换，分别舍弃坐标和奇异方向的幅度，实际位移、原目标残差和矩阵运算成本要分开核对。

自适应方法用历史幅度决定尺度。AdaGrad累计平方而难以遗忘，RMSProp用指数平均，Adam再加上方向平均和零初始化修正，AMSGrad让分母单调，AdamW把收缩与统计分开，Adafactor和Shampoo处理矩阵参数的存储和方向关系。这些尺度通常依赖当前读数，每种保证都必须连同它的假设一起引用，单调性之类的代数性质不是收敛结论。

裁剪、近似和替代让实际方向偏离梯度。把读数分解为真梯度、条件偏差和噪声，就能看到噪声项随步长缩小，偏差项不随步长缩小，而上界中的残余不一定是实际的下限。替代方向不是原目标的导数，评价和停止都要回到实际的前向目标。

所有保证都对应一个具体的目标、一个返回点和一个计算单位。同一个问题按更新次数和按读取量计时，方法的排序可以不同。有限步优化误差、路径的长期收敛和样本造成的统计误差是三个不同的问题，本章处理的是第一个，只说明了另外两个与它的分工。样本替换的分析则把算法看作从数据到输出的映射，为后续讨论统计误差提供了一种输入。

随机梯度方法需要目标的梯度读数。若只能读到函数值，就要先由函数值构造方向，下一章研究这个问题。
